\documentclass[12pt]{article}
\usepackage[a4paper,margin=1in]{geometry}
\usepackage[utf8]{inputenc}
\usepackage[T1]{fontenc}
\usepackage{amsmath,amssymb}
\usepackage{textcomp}
\usepackage{parskip}
\begin{document}
\begin{center}
{\Large \textbf{2011 STEP III}}
\end{center}

\subsection*{Page 1}
\begin{flushleft}
\bigskip\textbf{Section A: Pure Mathematics}\par
\bigskip\textbf{Question 1.} (i) Find the general solution of the differential equation\par
du x + 2\par
\[u = 0 .\]
dx - \lambda x + 1 \mu\par
\textbf{(ii)} Show that substituting y = ze\par
-x\par
(where z is a function of x) into the second order\par
differential equation\par
d y dy\par
\[(x + 1) + x y = 0 ( )\]
\[dx2 dx - ∗\]
dz\par
leads to a first order differential equation for . Find z and hence show that the\par
dx\par
general solution of ( ) is\par
∗\par
\[y = Ax + Be\]
-x\par
,\par
where A and B are arbitrary constants.\par
\textbf{(iii)} Find the general solution of the differential equation\par
d y dy\par
\[(x + 1) + x y = (x + 1) .\]
\[dx2 dx -\]
\bigskip\textbf{Question 2.} The polynomial f(x) is defined by\par
\[f(x) = xn + a xn-1 + + a x 2 + a x + a ,\]
n-1\par
\cdots\par
2 1 0\par
\[where n > 2 and the coefficients a\]
, . . . , a\par
n-1\par
\[are integers, with a^0\]
\[= 0. Suppose that the\]
equation f(x) = 0 has a rational root p/q, where p and q are integers with no common factor\par
\[greater than 1, and q > 0. By considering\]
qn-1\par
f(p/q), find the value of q and deduce that any\par
\[rational root of the equation f(x) = 0 must be an integer.\]
\textbf{(i)} Show that the nth root of 2 is irrational for n > 2.\par
\textbf{(ii)} Show that the cubic equation\par
\[x x + 1 = 0\]
-\par
has no rational roots.\par
\textbf{(iii)} Show that the polynomial equation\par
\[xn 5x + 7 = 0\]
-\par
\[has no rational roots for n > 2.\]
\[9475 Jun11\]
\end{flushleft}
\subsection*{Page 2}
\begin{flushleft}
2 3\par
\bigskip\textbf{Question 3.} Show that, provided q = 4p , the polynomial\par
\[x 3px + q (p = 0, q = 0)\]
- 6 6\par
can be written in the form\par
3 3\par
a(x α) + b(x β) ,\par
- -\par
2 2\par
\[where α and β are the roots of the quadratic equation pt qt + p = 0, and a and b are\]
-\par
constants which you should express in terms of α and β.\par
\[Hence show that one solution of the equation x 24x + 48 = 0 is\]
-\par
2(2 2 3 )\par
\[x = -\]
1 2 3\par
-\par
\[and obtain similar expressions for the other two solutions in terms of ω, where ω = e\]
\[2\pii/3\]
.\par
3 2 3\par
\[Find also the roots of x 3px + q = 0 when p = r and q = 2r for some non-zero constant r.\]
-\par
\bigskip\textbf{Question 4.} The following result applies to any function f which is continuous, has positive gradient and\par
\[satisfies f(0) = 0 :\]
a b\par
ab 6 f(x) dx + f\par
-1\par
(y) dy , ( )\par
\[Z^0 Z^0 ∗\]
\[where f -1 denotes the inverse function of f, and a > 0 and b > 0.\]
\textbf{(i)} By considering the graph of y = f(x), explain briefly why the inequality ( ) holds.\par
∗\par
\[In the case a > 0 and b > 0, state a condition on a and b under which equality holds.\]
1 1\par
\textbf{(ii)} By taking f(x) =\par
xp-1\par
\[in ( ), where p > 1, show that if + = 1 then\]
∗ p q\par
ap bq\par
ab 6 + .\par
p q\par
Verify that equality holds under the condition you stated above.\par
\textbf{(iii)} Show that, for 0 6 a 6 1 \pi and 0 6 b 6 1,\par
\[ab 6 b arcsin b + 1 b^2 cos a .\]
- -\par
p\par
\[Deduce that, for t > 1,\]
\[arcsin(t -1 ) > t t^2 1 .\]
- -\par
p\par
\[9475 Jun11 [Turn over\]
\end{flushleft}
\subsection*{Page 3}
\begin{flushleft}
\bigskip\textbf{Question 5.} A movable point P has cartesian coordinates (x, y), where x and y are functions of t. The\par
polar coordinates of P with respect to the origin O are r and θ. Starting with the expression\par
1 2\par
r dθ\par
\[\int\]
for the area swept out by OP , obtain the equivalent expression\par
dy dx\par
x y dt . ( )\par
Z \lambda dt - dt \mu ∗\par
The ends of a thin straight rod AB lie on a closed convex curve . The point P on the rod\par
C\par
is a fixed distance a from A and a fixed distance b from B. The angle between AB and the\par
positive x direction is t. As A and B move anticlockwise round , the angle t increases from\par
C\par
\[0 to 2\pi and P traces a closed convex curve inside , with the origin O lying inside , as\]
D C D\par
shown in the diagram.\par
y\par
B\par
b\par
P C\par
D\par
a\par
t x\par
O\par
A\par
Let (x, y) be the coordinates of P . Write down the coordinates of A and B in terms of a, b,\par
x, y and t.\par
The areas swept out by OA, OB and OP are denoted by [A], [B] and [P ], respectively. Show,\par
using ( ), that\par
∗ 2\par
\[[A] = [P ] + \pia af\]
-\par
where\par
\[2\pi\]
dy dx\par
\[f = x + cos t + y sin t dt .\]
\[Z^0 \lambda \gamma dt \delta \gamma - dt \delta \mu\]
Obtain a corresponding expression for [B] involving b. Hence show that the area between the\par
\[curves and is \piab.\]
C D\par
\[9475 Jun11\]
\end{flushleft}
\subsection*{Page 4}
\begin{flushleft}
\bigskip\textbf{Question 6.} The definite integrals T , U , V and X are defined by\par
1 ln 3\par
2 artanh t u\par
\[T = dt , U = du ,\]
Z 1 t Zln 2 2 sinh u\par
1 1 ln 3\par
2 ln v 2\par
\[V = dv , X = ln(coth x) dx .\]
\[- Z 1 1 v^2 Z 1 ln 2\]
3 - 2\par
Show, without evaluating any of them, that T , U , V and X are all equal.\par
7 Let\par
n\par
\[T = \sqrt\{a\} + 1 + \sqrt\{a\} ,\]
n\par
where n is a positive integer and a is any given positive integer.\par
\textbf{(i)} In the case when n is even, show by induction that T can be written in the form\par
n\par
A + B a(a + 1) ,\par
n n\par
p\par
2 2\par
\[where A and B are integers (depending on a and n) and A = a(a + 1)B + 1.\]
n n n n\par
\textbf{(ii)} In the case when n is odd, show by considering (\sqrta + 1 + \sqrta)T where m is even, or\par
m\par
otherwise, that T can be written in the form\par
n\par
\[C \sqrt\{a\} + 1 + D \sqrt\{a\} ,\]
n n\par
2 2\par
\[where C and D are integers (depending on a and n) and (a + 1)C = aD + 1 .\]
n n n n\par
\textbf{(iii)} Deduce that, for each n, T can be written as the sum of the square roots of two\par
n\par
consecutive integers.\par
\bigskip\textbf{Question 8.} The complex numbers z and w are related by\par
1 + iz\par
\[w = .\]
i + z\par
\[Let z = x + iy and w = u + iv, where x, y, u and v are real. Express u and v in terms of x\]
and y.\par
\textbf{(i)} By setting x = tan(θ/2), or otherwise, show that if the locus of z is the real axis y = 0,\par
2 2\par
\[< x < , then the locus of w is the circle u + v = 1 with one point omitted.\]
-\infty \infty\par
\textbf{(ii)} Find the locus of w when the locus of z is the line segment y = 0, 1 < x < 1 .\par
-\par
\textbf{(iii)} Find the locus of w when the locus of z is the line segment x = 0, 1 < y < 1 .\par
-\par
\textbf{(iv)} Find the locus of w when the locus of z is the line y = 1, < x < .\par
-\infty \infty\par
\[9475 Jun11 [Turn over\]
\end{flushleft}
\subsection*{Page 5}
\begin{flushleft}
\bigskip\textbf{Section B: Mechanics}\par
9 Particles P and Q have masses 3m and 4m, respectively. They lie on the outer curved surface\par
of a smooth circular cylinder of radius a which is fixed with its axis horizontal. They are\par
connected by a light inextensible string of length \pia, which passes over the surface of the\par
cylinder. The particles and the string all lie in a vertical plane perpendicular to the axis of\par
the cylinder, and the axis intersects this plane at O. Initially, the particles are in equilibrium.\par
Equilibrium is slightly disturbed and Q begins to move downwards. Show that while the two\par
particles are still in contact with the cylinder the angle θ between OQ and the vertical satisfies\par
7aθ˙2\par
\[+ 8g cos θ + 6g sin θ = 10g .\]
\textbf{(i)} Given that Q loses contact with the cylinder first, show that it does so when θ = β,\par
where β satisfies\par
\[15 cos β + 6 sin β = 10.\]
\textbf{(ii)} Show also that while P and Q are still in contact with the cylinder the tension in the\par
string is mg(sin θ + cos θ) .\par
10 Particles P and Q, each of mass m, lie initially at rest a distance a apart on a smooth\par
horizontal plane. They are connected by a light elastic string of natural length a and modulus\par
1 2\par
of elasticity maω , where ω is a constant.\par
Then P receives an impulse which gives it a velocity u directly away from Q. Show that when\par
the string next returns to length a, the particles have travelled a distance \piu/ω , and find\par
the speed of each particle.\par
Find also the total time between the impulse and the subsequent collision of the particles.\par
\[9475 Jun11\]
\end{flushleft}
\subsection*{Page 6}
\begin{flushleft}
\bigskip\textbf{Question 11.} A thin uniform circular disc of radius a and mass m is held in equilibrium in a horizontal\par
plane a distance b below a horizontal ceiling, where b > 2a. It is held in this way by n light\par
inextensible vertical strings, each of length b; one end of each string is attached to the edge\par
of the disc and the other end is attached to a point on the ceiling. The strings are equally\par
spaced around the edge of the disc. One of the strings is attached to the point P on the\par
disc which has coordinates (a, 0, b) with respect to cartesian axes with origin on the ceiling\par
-\par
directly above the centre of the disc.\par
The disc is then rotated through an angle θ (where θ < \pi) about its vertical axis of symmetry\par
and held at rest by a couple acting in the plane of the disc. Show that the string attached\par
to P now makes an angle φ with the vertical, where\par
\[b sin φ = 2a sin θ .\]
Show further that the magnitude of the couple is\par
mga sin θ\par
.\par
\[b^2 4a2 sin 2 1 θ\]
q -\par
The disc is now released from rest. Show that its angular speed, ω, when the strings are\par
vertical is given by\par
2 2\par
a ω\par
\[= b b^2 4a2 sin 2 1 θ .\]
4g - q -\par
\[9475 Jun11 [Turn over\]
\end{flushleft}
\subsection*{Page 7}
\begin{flushleft}
\bigskip\textbf{Section C: Probability and Statistics}\par
\bigskip\textbf{Question 12.} The random variable N takes positive integer values and has pgf (probability generating\par
\[function) G(t). The random variables X , where i = 1, 2, 3, . . . , are independently and\]
i\par
identically distributed, each with pgf H(t). The random variables X are also independent of\par
i\par
N . The random variable Y is defined by\par
N\par
\[Y = X .\]
i\par
\[Xi=1\]
Given that the pgf of Y is G(H(t)), show that\par
\[E(Y ) = E(N )E(X ) and Var (Y ) = Var (N ) E(X ) + E(N )Var (X ) .\]
i i i\par
A fair coin is tossed until a head occurs. The total number of tosses is N . The coin is then\par
tossed a further N times and the total number of heads in these N tosses is Y . Find in this\par
\[particular case the pgf of Y , E(Y ), Var (Y ) and P(Y = r).\]
\bigskip\textbf{Question 13.} In this question, the notation x denotes the greatest integer less than or equal to x, so for\par
b c\par
\[example \pi = 3 and 3 = 3.\]
b c b c\par
\textbf{(i)} A bag contains n balls, of which b are black. A sample of k balls is drawn, one after\par
another, at random with replacement. The random variable X denotes the number of\par
black balls in the sample. By considering\par
\[P(X = r + 1)\]
,\par
\[P(X = r)\]
show that, in the case that it is unique, the most probable number of black balls in the\par
sample is\par
(k + 1)b\par
.\par
\[\phi n \pi\]
Under what circumstances is the answer not unique?\par
\textbf{(ii)} A bag contains n balls, of which b are black. A sample of k balls (where k 6 b) is drawn,\par
one after another, at random without replacement. Find, in the case that it is unique,\par
the most probable number of black balls in the sample.\par
Under what circumstances is the answer not unique?\par
\[9475 Jun11\]
\end{flushleft}
\end{document}
